%% The line mask, reported in two layers: the frequency coincidence, which is
%% a velocity criterion and nothing else, and the one species we do not read
%% as emission, which is an argument and is set out as one.
%% Owner: r16-mask.  Carries \label{app:linemask} (cited by 04_method.tex and
%% app_C_spatialnorm.tex) and \label{tab:masklines}.
%%
%% ★ THIS APPENDIX IS THE PAPER'S ONLY DEFINITION OF THE ATTRIBUTION RULE.
%% It owns the species list, the half-width and the FRAME.  The rule is the
%% star's own rest frame; the local standard of rest is evaluated as a check
%% and changes no disposition (masklsr_v520.py, L10).  Sec. 5.3 states the
%% rule in one clause and Appendix D points here; nothing else in the paper
%% may name a frame of its own, or there are two rules again.
\section{The line mask}
\label{app:linemask}

A crossing that falls on a known molecular transition has an astrophysical
explanation to hand, and the mask exists to say which crossings those are,
reporting the frequency coincidence and the judgement about the coinciding
species separately.

\emph{The list, and when it was made.} The mask holds every transition of the
\FrNSpec{} species these spectral setups cover -- the three CO isotopologues,
HCN, HCO$^{+}$, CS, CN, SiO, H$_2$CO, [C\,\textsc{i}] and hydrogen
recombination -- that falls inside one of the \McNIsland{} searched frequency
intervals with an upper-state energy below \McEuMax\,K. That rule is the whole
of it, and it admits \FrNTrans{} transitions (Table~\ref{tab:masklines}).
No strength cut is applied across species: any threshold admitting HCN's
strong lines would discard the brightest lines in these fields. Within a
species the main hyperfine components are kept and the satellites
dropped. The list is a single Splatalogue query
\citep{Splatalogue2007} of CDMS and JPL \citep{Pickett1998,CDMS2005} made on
\MkListDate{}, and the query is dated and repeatable, so the mask itself can
be rebuilt.

\emph{The species are the ones the observers went to the telescope for, and
that is a limitation.} The spectral setups are tuned to them: \TnCenPct{} per
cent of the \NWinA{} Class~A windows are centred within
$\pm\MaskHalfKms$\,km\,s$^{-1}$ of a transition of the list and \TnWinPct{}
per cent contain one somewhere, against \TnCtrlCenLo--\TnCtrlCenHi{} per cent
for the same windows displaced bodily. The list is therefore set by the
science goals of the programmes that took the data and not by anything in this
search, and \emph{unattributed} means unattributed against it. A wider list
measures what that costs: every transition below \XlEuMax\,K of all
\XlNSpecies{} species with a published interstellar or circumstellar detection
\citep{Pickett1998,CDMS2005} gives \XlNTrans{} transitions inside the searched
intervals, covering \XlBandPct{} per cent of the Class~A union against
\MkMaskAPct{} per cent for the list adopted here, and it attributes every one
of the \MkUnattr{} unattributed crossings together with \XlShareWidePct{} per
cent of the control positions' own exceedances, so it empties the category for
star and control alike; narrowed to $\pm\XlTightHalf$\,km\,s$^{-1}$ it would
attribute \XlTightNIn{} of the \MkUnattr{} and the unattributed count remains
what the control positions predict, \XlTightObs{} against \XlTightExp{}
expected. The nearest transition of that wider list is reported beside every
unattributed crossing in the accompanying data -- and at a density where
\XlBandPct{} per cent of the band lies within \MaskHalfKms\,km\,s$^{-1}$ of
something, proximity is not an identification.

\emph{What an identification would look like, on the pair that comes closest
to one.} The \ZmNCross{} unattributed crossings of $\beta$~Pictoris lie
\ZmOffLoMHz{} and \ZmOffHiMHz\,MHz above components of the methanol
$5_K\to4_K$ ladder (\ZmCompLo{} and \ZmCompHi, $E_{\rm u}=\ZmEuLo$ and
\ZmEuHi\,K, at frequencies on which CDMS and JPL agree), \ZmDvLo{} and
\ZmDvHi\,km\,s$^{-1}$ in the star's frame: offsets agreeing with each other
to \ZmDvAgree\,km\,s$^{-1}$, which is the configuration a real line
identification produces and which arises by chance in \ZmChancePct{} per cent
of pairs drawn uniformly in the window they were found in. Four measurements
decide against it, all in the one window that covers the whole ladder in both
blocks, \ZmBandLo--\ZmBandHi\,GHz in the stellar frame. The brightest member of the complex, \ZmBright{} at
\ZmBrightMHz\,MHz, would reach $T_\star\ge\ZmPredAplus$ at any excitation
temperature and would therefore have been the window maximum, where the
maximum over the whole channel--drift grid is the crossing itself. The weaker
of the two claimed components cannot exceed \ZmRatioMax{} of the stronger at
any temperature and is observed at \ZmRatioObs{} of it. Both fitted drifts
are \ZmStepsOne{} and \ZmStepsTwo{} steps of their own grid from the zero a
transition at rest requires, each track crossing \ZmSweepOne{} and
\ZmSweepTwo{} channels inside its own block. And the two are not in one
block: the first is in \texttt{\ZmBlockOne} and the second in
\texttt{\ZmBlockTwo}, \ZmGapH\,h apart, and each is absent from the other,
where a persistent line would have been recovered at \ZmPersOne{} and
\ZmPersTwo$\sigma$ and the statistic at that cell reaches only \ZmAnyOne{}
and \ZmAnyTwo{} over every drift of the grid (\S\ref{sec:etacrv}). A real
pair of components would have stood above the trigger in both blocks; one in
each and neither in both is what a noise pair gives. Where this system's own molecular gas
is seen, in \ZmCoNBlock{} blocks, it registers to \ZmCoReg\,km\,s$^{-1}$ of
rest.

\emph{Attribution is decided in the star's own rest frame, and the sign is
not the usual one.} A crossing's frequency is carried from the topocentric
value to the barycentre using its own block's pointing and mid-time, and then
to the stellar rest frame using the star's systemic velocity. That frame, and
no other, is where a coincidence is judged, a tolerance applied in the
observed frame spending most of its budget on the Earth's motion. The
\FrEpochNEpoch{} \FrEpochStar{} epochs show the transform working in the only
way that can be checked: \FrEpochTopoKms\,km\,s$^{-1}$ apart as observed, they
agree to \FrEpochStelKms\,km\,s$^{-1}$ in the star's frame.
Throughout, $\Delta v_\star = c\,(\nu_\star-\nu_0)/\nu_0$, so a positive
offset means the crossing lies \emph{above} the transition
in frequency, the opposite sign to the velocity convention of a spectral line.

\emph{Every crossing is read in the local standard of rest as well, where
interstellar gas is at rest, and that second reading
establishes which coincidences are interstellar.} Of the \LsNLos{} stars
carrying crossings, the composite CO($1\to0$) survey of
\citet{DameThaddeus2022,Dame2001} covers \LsNLosCov{} and holds emission
toward \LsNLosDet{} of them, while toward every other star, including the
\LsNLosOut{} outside its footprint, the all-sky Planck CO map
\citep{PlanckCO2016} stays below \LsPlanckMax\,K\,km\,s$^{-1}$. Toward
\LsCloudStars{} the CO spectrum is non-zero at the crossing's own LSR
velocity, within \LsCloudDvMax\,km\,s$^{-1}$ of its own peak and at
\LsCloudTMin\,K or brighter against a \LsCoRms\,K channel, so \LsNCloud{} of
the \MkNCoinc{} coincidences are interstellar gas and not circumstellar.
Two are worth naming. \BkAlmaName{} is a serendipitous field star toward the
Lupus clouds, with the clouds far behind it, and its CO($2\to1$) crossing lies
\LsLupDv\,km\,s$^{-1}$ from a \LsLupWco\,K\,km\,s$^{-1}$ cloud line at
\LsLupVco\,km\,s$^{-1}$: a molecular sightline and not a disc, on which the
wider chemistry a cloud carries is a real limitation, \XlLupNamedN{} of the
species such a sightline would show reaching one of its crossings within the
mask's own half-width, \XlLupNamedMol{} at \XlLupNamedDv\,km\,s$^{-1}$.
HD~48370 is the case where the two readings part: its coincidence lies
\LsHdDvStar\,km\,s$^{-1}$ from the star's own CO($2\to1$) while in the local
standard of rest it sits at \LsHdVfeat\,km\,s$^{-1}$, on a catalogued cloud
line peaking at \LsHdVco, so the offset follows the interstellar gas and not
the star and that window's extended emission is a cloud and not a resolved
disc.

Table~\ref{tab:masklines} gives the
mask species by species; the hydrogen recombination lines among them carry
no tabulated upper-state energy, and are blank in that column.

\begin{table}
\centering
\footnotesize
%% ★ v4.26 (r25-captions): 4 printed lines -> the author's 3.  The hydrogen
%% recombination lines' missing $E_{\rm u}$ MOVED to the paragraph above: it
%% explains a blank cell, which is a fact about the catalogue and not a
%% definition of a column.
\caption{The line mask, by species: the number $N$ of transitions, the
frequency range they span inside the searched intervals and their
upper-state energies.}
\label{tab:masklines}
%% GENERATED by maskso_v414.py -- do not hand-edit.
\begin{tabular}{@{}lrr@{\,}c@{\,}rr@{}}
\hline
Species & $N$ & \multicolumn{3}{c}{$\nu_0$ (GHz)} & $E_{\rm u}$ (K) \\
\hline
H & 5 & 92.03 & -- & 231.90 & -- \\
H$_2$CO & 13 & 101.33 & -- & 491.97 & 10--106 \\
C$^{18}$O & 3 & 109.78 & -- & 329.33 & 5--32 \\
CN & 50 & 113.12 & -- & 680.30 & 5--114 \\
CO & 3 & 115.27 & -- & 345.80 & 6--33 \\
CS & 3 & 195.95 & -- & 342.88 & 24--66 \\
SiO & 2 & 217.10 & -- & 347.33 & 31--75 \\
$^{13}$CO & 8 & 220.40 & -- & 330.59 & 16--32 \\
HCN & 7 & 354.50 & -- & 354.51 & 43--43 \\
HCO$^{+}$ & 1 & 356.73 & -- & 356.73 & 43 \\
{[C\,\textsc{i}]} & 1 & 492.16 & -- & 492.16 & 24 \\
\hline
\end{tabular}%

\end{table}

\emph{The coincidences, and the one species the mask leaves out.} Of the
\NCross{} crossings, \MkNCoinc{} fall within $\pm\MaskHalfKms$\,km\,s$^{-1}$
of a transition of the list in their own star's frame. The list holds no
species whose coincidences are reported and not acted on, so those are also
the \MkAttr{} attributions and the two counts are one count. Sulphur monoxide
is a catalogued species of these spectral setups and is deliberately not in
the list -- no debris disc has a reported SO detection, where the
second-generation gas is CO, C and O, so masking its \SoNTrans{} transitions
would have declared no constraint over bandwidth this survey does constrain --
and the \SoNCrossWord{} crossings within the same half-width of one of them
(\SoStars, at \SoDvList\,km\,s$^{-1}$ from transitions of upper-state energy
\SoEuLo--\SoEuHi\,K) are therefore searched, tested and reported as
unattributed, with the offset printed beside each in the ledger and
\SoNBelowScreen{} of them below the rank screen as well. Interstellar sulphur
monoxide is a measurement rather than a preference: in the local standard of
rest the three offsets are \SoDvLsrList\,km\,s$^{-1}$, so for one of them the
coincidence does not survive the change of frame at all, \SoNUncovered{} lie
outside the composite survey's footprint on sightlines carrying at most
\SoPlanckMax\,K\,km\,s$^{-1}$ of integrated CO against the
\SoPlanckCeil\,K\,km\,s$^{-1}$ a detection needs, and the third,
\SoCloudStar, is a molecular sightline whose crossing nevertheless sits
\SoCloudDv\,km\,s$^{-1}$ from the emission, where this survey's own cloud test
allows \LsCloudDvMax\,km\,s$^{-1}$. Each of the three windows also covers
CO($3\to2$) and in each the maximum over the whole channel--drift grid is the
crossing itself, so no CO($3\to2$) feature above the trigger accompanies any
of them -- and carbon monoxide, not sulphur monoxide, is what a cloud shows
first. The reserved hold-out is treated identically, leaving \HoNAttr{}
attributed and \HoNUnattr{} unattributed there
(\S\ref{sec:holdoutsearch}).

\emph{The half-width, and what it costs.} A one-channel test is too tight for
circumstellar gas, whose emission is shifted by many channels: the width
follows from Keplerian motion where these molecules survive,
$\lesssim$13\,km\,s$^{-1}$ beyond 10\,au, and is
$\pm\MaskHalfKms$\,km\,s$^{-1}$. At that width the mask removes
\MkMaskAGHz\,GHz of the \CvUnionA\,GHz Class~A union, \MkMaskAPct{} per cent
of it, and \FrMaskLostGHz\,GHz of \FrMaskUnionGHz\,GHz over both classes
together, \FrMaskPct{} per cent, so the cost falls on the primary
experiment, because its windows are the line-tuned ones.
Attributing on either frame instead would mask the union of two tubes
separated by the star's LSR velocity, \LsMaskTwoGHz\,GHz or
\LsMaskTwoPct{} per cent of the Class~A union -- \LsMaskAddPct{} per cent
more bandwidth for no further attribution, which is the measurement behind
deciding in one frame.

\emph{The labelled crossings are tested for recurrence too}
(\S\ref{sec:etacrv}). For the \MkAttrExcl{} not recovered a repeat window's
own largest statistic bounds the predicted cell from above and the ledger
prints the row as a bound, at $T_{\rm rep}<\RqHrTRep$ and
$T_{\rm rep}<\RqAjTRep$ against the $T_{\rm pers}=\RqHrTPers$ and
$T_{\rm pers}=\RqAjTPers$ real carbon monoxide at the discovery flux density
would have given. One of those two attributions rests on a velocity rather
than on the recurrence test, the sightline carrying catalogued interstellar
emission \LsLupDv\,km\,s$^{-1}$ from the crossing's own LSR velocity, and its
repeat blocks show where that gas lies rather than whether it can be
recovered: in every one of them between \RqAjNgeLo{} and \RqAjNgeHi{} of the
\NCtrl{} control positions stand above the star, so the gas is bright and it
is not at the star.
